\chapter{Geometric Inevitability --- From Numerology to Derivation}
\label{app:geometric_inevitability}

\textit{Every ``mystical'' constant in physics is a geometric packing theorem in disguise.}

\vspace{1em}

Throughout the history of mathematics and natural philosophy, certain constants and sequences have been observed repeatedly across wildly disparate physical systems. The Golden Ratio ($\varphi$), the Fibonacci sequence, $\pi$, and the nuclear ``magic numbers'' have each spawned vast literatures of numerological speculation---claims that these quantities reflect some deep, unexplained cosmic design principle.

The Applied Vacuum Engineering framework resolves each of these cases. None require mysticism. All are geometric inevitabilities: forced outcomes of minimum-impedance packing constraints applied to finite numbers of coupled resonant nodes in 3D space. This appendix traces each constant from its traditional numerological invocation to its deterministic derivation.

\section{The Golden Ratio: Minimum Impedance at 12 Nodes}
\label{sec:golden_ratio_emergence}

The Golden Ratio,
\begin{equation}
    \varphi = \frac{1 + \sqrt{5}}{2} \approx 1.61803\ldots
\end{equation}
appears throughout nature: sunflower seed spirals, nautilus shell curvatures, DNA helical geometry, galactic arm spacing. These appearances have been interpreted as evidence of mathematical design principles transcending physical law.

\subsection{Derivation}
The regular icosahedron is the unique Platonic solid with 12 vertices. Its vertex coordinates are permutations of:
\begin{equation}
    (0, \pm 1, \pm \varphi)
\end{equation}
This is not a choice; it is the \textit{only} solution to the Thomson problem (distributing 12 identical repulsive charges on a sphere to minimize total potential energy). No other arrangement of 12 equidistant points on $S^2$ exists.

When the AVE semiconductor binding engine solves Chromium-52 ($Z=24$, $A=52$) as 13 alpha clusters (12 on an icosahedral shell + 1 at the center), the solver converges to:
\begin{equation}
    R_{\text{ico}} = 166.5\,d, \quad \text{Error} = 0.0001\%
\end{equation}
using only the axiom-derived coupling constant $K = (5\pi/2) \cdot \alpha\hbar c / (1 - \alpha/3) \approx 11.337$ MeV$\cdot$fm and zero empirical parameters.

\subsection{Physical Interpretation}
$\varphi$ does not ``choose to appear'' in Chromium-52. It is \textbf{forced into existence} by the requirement that 12 mutually repulsive inductive loads (alpha clusters) minimize their total reactive impedance on a bounded spherical surface. The Thomson problem has exactly one solution at $N=12$: the icosahedron. The icosahedron is defined by $\varphi$. Therefore, any physical system that must distribute 12 coupled nodes on a sphere will exhibit $\varphi$---whether it is a nuclear shell, a viral capsid, or a fullerene cage.

The ubiquity of $\varphi$ in nature is not mysterious. It is the inevitable geometric consequence of packing 12 objects symmetrically in 3D.

\section{The Fibonacci Sequence: Convergent Ratio as Packing Proxy}
\label{sec:fibonacci_packing}

The Fibonacci sequence ($1, 1, 2, 3, 5, 8, 13, 21, 34, \ldots$) exhibits the property that the ratio of consecutive terms converges to $\varphi$:
\begin{equation}
    \lim_{n \to \infty} \frac{F_{n+1}}{F_n} = \varphi
\end{equation}

In the AVE periodic table catalog, elements beyond $Z=28$ (where no exact geometry has been solved) use a \textit{Fibonacci lattice} as a numerical proxy for distributing alpha clusters on a sphere. This is not coincidental: the Fibonacci lattice is the computationally cheapest algorithm that produces approximately uniform point distributions on $S^2$---because it converges to the icosahedral ground state that defines $\varphi$.

The fact that the Fibonacci proxy achieves $<0.5\%$ mass prediction accuracy across 105 elements (Volume VI, Appendix A) is itself a proof that the underlying physical packing is icosahedral in character. The Fibonacci lattice works \textit{because} it approximates the true minimum-impedance geometry.

\section{Pi and the Topological Horizon}
\label{sec:pi_horizon}

The constant $\pi$ appears in the AVE framework at a highly specific structural boundary: the \textbf{Topological Horizon} of Boron-11.

When solving the 7-nucleon halo distance for Boron-11 ($Z=5$, $A=11$), the semiconductor engine places the halo at:
\begin{equation}
    R_{\text{halo}} = 11.84\,d
\end{equation}
The theoretical maximum distance before topological decoupling is the full isotropic solid angle boundary:
\begin{equation}
    \text{Horizon}_{\text{limit}} = 4\pi - \frac{\sqrt{2}}{2} \approx 11.859
\end{equation}
The proximity of $R_{\text{halo}}$ to $4\pi - \sqrt{2}/2$ means that Boron-11's halo sits at the absolute razor edge of the topological decoupling boundary. The factor $4\pi$ arises because $4\pi$ steradians is the total solid angle of a sphere---the complete isotropic radiation boundary of a point source.

$\pi$ does not appear here because of some universal principle of harmony. It appears because Boron-11 is a spherical source radiating into 3D Euclidean space, and the total solid angle of a sphere is exactly $4\pi$. This is Gauss's law, not numerology.

\section{Nuclear Magic Numbers: Shell Closure as Impedance Matching}
\label{sec:magic_numbers}

Standard nuclear physics identifies the ``magic numbers'' ($2, 8, 20, 28, 50, 82, 126$) as nucleon counts at which nuclei are anomalously stable. These numbers are traditionally derived from the nuclear shell model with spin-orbit coupling---an empirical fit to observed binding energies.

In the AVE framework, each magic number corresponds to a geometry where the alpha-cluster packing achieves impedance matching ($S_{11} \to 0$, maximum Q-factor):
\begin{center}
\renewcommand{\arraystretch}{1.3}
\begin{tabular}{r l l l}
\toprule
\textbf{Magic $Z$} & \textbf{$N_\alpha$} & \textbf{Geometry} & \textbf{Structural Meaning} \\
\midrule
$2$  & $1\alpha$  & Single Tank       & Borromean braid closure \\
$8$  & $4\alpha$  & Tetrahedron       & Minimum 3D Platonic solid \\
$20$ & $10\alpha$ & Bicapped Antiprism & Maximum pre-Archimedean packing \\
$28$ & $14\alpha$ & FCC-14            & Face-centered cubic closure \\
\bottomrule
\end{tabular}
\end{center}

The underlying pattern is clean geometric closure: a complete Platonic or Archimedean solid where every alpha cluster is symmetrically equivalent and the total strain field has zero net dipole moment.

\section{The Platonic Progression}

The systematic walk through nuclear topologies traces an exact path through the classical Platonic and Archimedean solids:

\begin{center}
\renewcommand{\arraystretch}{1.3}
\begin{tabular}{r l l r r}
\toprule
\textbf{$N_\alpha$} & \textbf{Element} & \textbf{Geometry} & \textbf{$R$ ($d$)} & \textbf{Error} \\
\midrule
$3$  & C-12  & Triangle (Ring)        & 56.5  & 0.000000\% \\
$4$  & O-16  & Tetrahedron            & 33.4  & 0.000000\% \\
$5$  & Ne-20 & Pentagonal Ring        & 81.2  & 0.000000\% \\
$6$  & Mg-24 & Octahedron             & 78.0  & 0.000000\% \\
$7$  & Si-28 & Pentagonal Bipyramid   & 83.0  & 0.0002\%   \\
$8$  & S-32  & Cube                   & 4.7   & 0.000000\% \\
$10$ & Ar-40 & Bicapped Antiprism     & 94.0  & 0.0002\%   \\
$10$ & Ca-40 & Bicapped Antiprism     & 5.9   & 0.000000\% \\
$12$ & Ti-48 & Cuboctahedron          & 116.5 & 0.0001\%   \\
$13$ & Cr-52 & Centered Icosahedron   & 166.5 & 0.0001\%   \\
$14$ & Fe-56 & FCC-14                 & 186.5 & 0.0001\%   \\
\bottomrule
\end{tabular}
\end{center}

This progression is not a fit. It is the unique sequence of minimum-impedance packing solutions for increasing numbers of mutually repulsive resonant nodes on a sphere. Each geometry is forced by the topology; the engine simply identifies which Platonic or Archimedean solid satisfies the impedance matching condition.

\section{Derived Numerical Constants}
\label{sec:derived_numerology}

Beyond the classical ``magic numbers'' of nuclear physics, the AVE
framework derives several numerical constants that govern phenomena
across vastly different scales.  Each is a ratio, product, or
geometric projection of two or more independently derived quantities.
None is fitted; all emerge from the axioms.

{
\renewcommand{\arraystretch}{1.35}
\small
\begin{longtable}{@{}lllp{3.8cm}@{}}
\caption{Derived Numerical Constants} \\
\toprule
\textbf{Constant} & \textbf{Formula} & \textbf{Value} & \textbf{Physical Meaning} \\
\midrule
\endfirsthead
\caption[]{Derived Numerical Constants (\textit{Continued})} \\
\toprule
\textbf{Constant} & \textbf{Formula} & \textbf{Value} & \textbf{Physical Meaning} \\
\midrule
\endhead
\bottomrule
\endfoot
$\delta_{th}$ & $\dfrac{1}{14\pi^2}
    = \dfrac{\nu_{vac}}{\kappa_{cold}} \cdot \dfrac{2}{\pi}$
    & 0.00724
    & Residual thermal softening \\[6pt]
$K_{mut}$ & $\dfrac{5\pi}{2}\cdot\dfrac{\alpha\hbar c}{1-\alpha/3}$
    & 11.34 MeV$\cdot$fm
    & Nuclear mutual inductance \\[6pt]
$d_p$ & $\dfrac{4\hbar}{m_p c}$
    & 0.841 fm
    & Proton charge radius \\[6pt]
$a_0$ & $\dfrac{c \cdot H_\infty}{2\pi}$
    & $1.07\times10^{-10}$ m/s$^2$
    & MOND acceleration scale \\[6pt]
$V_{halo}$ & Skew-line integral of Borromean link
    & 2.0
    & Toroidal crossing volume \\[6pt]
\midrule
\multicolumn{4}{c}{\textit{New: from cross-scale gap closure}} \\[4pt]
$V_{GW}/V_{snap}$ & $\dfrac{h \cdot c \cdot \ell_{node} \cdot 2\pi f}{V_{snap}}$
    & $1.42\times10^{-28}$
    & GW linearity ratio ($h=10^{-21}$) \\[6pt]
$\Gamma_{tach}$ & $\dfrac{Z_{rad} - Z_{conv}}{Z_{rad} + Z_{conv}}$
    & 0.818 $\approx$ 9/11
    & Tachocline reflection \\[6pt]
$r_s(\odot)$ & $\dfrac{2GM_\odot}{c^2}$
    & 2954 m
    & Solar Schwarzschild radius \\[6pt]
$\beta_\text{fold}$ & $\ln(3) \times 3/7$
    & 0.471
    & Protein folding barrier exponent \\[6pt]
$E_{yield}$ & $\dfrac{V_{yield}}{\ell_{node}}
    = \dfrac{\sqrt{\alpha}\, m_e^2 c^3}{e\, \hbar}$
    & $1.13\times10^{17}$ V/m
    & Dielectric yield field ($\sqrt{\alpha}\times E_S$) \\[6pt]
\midrule
\multicolumn{4}{c}{\textit{New: from packing geometry (Axioms~1, 2)}} \\[4pt]
$\varphi$ & $\dfrac{\pi\sqrt{2}}{6}$
    & 0.7405
    & FCC lattice packing fraction \\[6pt]
$1 - \varphi$ & $1 - \dfrac{\pi\sqrt{2}}{6}$
    & 0.2595
    & Void fraction (Op10 drain floor) \\[6pt]
$Y_c$ & $1 - \sqrt{1 - \varphi}$
    & 0.4906
    & Drain $\to$ floor transition \\[6pt]
\midrule
\multicolumn{4}{c}{\textit{New: from cooperative lattice yield (Axioms~1, 2, 4)}} \\[4pt]
$n_{coop}$ & $6_{edges} \times \dfrac{D}{D-1} = 7 \times (1+\nu)$
    & 9
    & Cooperative phase amplification \\[6pt]
$E_{HB}$ & $U_{\text{Op4,raw}} \times (1 - \varphi)$
    & 0.2158 eV
    & Hydrogen bond energy \\[6pt]
$T_c^{\text{water}}$ & $\dfrac{E_{HB}}{n_{coop} \cdot k_B}$
    & 278.3 K (5.1$^\circ$C)
    & Cooperative yield temperature \\[6pt]
\midrule
\multicolumn{4}{c}{\textit{New: from proton transfer eigenmode (Axioms~1--4)}} \\[4pt]
$k_{hb}$ & $\dfrac{2 E_{HB}}{d_{hb}^2}$
    & 2.248 N/m
    & H-bond spring constant (Op4 curvature) \\[6pt]
$k_{OH}$ & Coulomb bond solver
    & 791 N/m
    & O--H covalent spring constant \\[6pt]
$T_m$ & $\dfrac{\hbar}{k_B}\sqrt{\dfrac{k_{hb}}{m_H}}$
    & 279.5 K ($+6.3^\circ$C)
    & Water melting eigenmode ($+2.3\%$) \\[6pt]
$\theta_{HOH}$ & $\arccos\!\left(\dfrac{-1}{3(1+\Gamma)}\right)$
    & $104.48^\circ$
    & H-O-H bond angle ($+0.03^\circ$) \\[6pt]
$V_{II}$ & $V_I \times \varphi_{FCC}$
    & $23.11\ \text{\AA}^3$
    & Topological cell collapse (State II) \\[6pt]
\end{longtable}
}

\subsection{Thermal Softening ($\delta_{th} = 1/(14\pi^2)$)}

The proton's core temperature $\sim m_p c^2 / k_B \approx 10^{13}$~K
softens the quartic Skyrme repulsion by partially averaging out the
gradient tensor.  With the lattice gradient saturation (Axiom~4) now
handled inside the energy functional, the residual softening fraction
is the RMS noise averaging:
\begin{equation}
    \delta_{th} = \frac{\nu_{vac}}{\kappa_{cold}} \cdot \frac{2}{\pi}
    = \frac{2/7}{8\pi} \cdot \frac{2}{\pi} = \frac{1}{14\pi^2} \approx 0.00724
\end{equation}
Both $\nu_{vac}$ and $\kappa_{cold}$ are pure geometric constants
derived independently; the $2/\pi$ factor is the mean-to-peak ratio
of rectified sinusoidal thermal noise.

\subsection{Nuclear Mutual Inductance ($K_{mut}$)}

The pairwise binding between nucleons is governed by:
\begin{equation}
    K_{mut} = \frac{5\pi}{2} \cdot \frac{\alpha \hbar c}{1 - \alpha/3}
    \approx 11.337\;\text{MeV}\cdot\text{fm}
\end{equation}
The factor $5\pi/2$ arises from five crossings of the proton's
$(2,5)$ cinquefoil knot, each contributing a $\pi/2$ phase advance
to the flux-linkage integral.  The $(1 - \alpha/3)^{-1}$ correction
is the first-order proximity enhancement at nuclear separations,
analogous to a transformer vertex correction.

\subsection{Proton Charge Radius ($d_p = 4\hbar/(m_p c)$)}

The proton charge radius is four times the proton reduced Compton
wavelength --- the RMS vibration amplitude of the centre-of-mass
standing wave confined within the saturated ($0\,\Omega$) cavity
boundary of one lattice cell:
\begin{equation}
    d_p = 4 \cdot \frac{\hbar}{m_p \, c} \approx 0.841\;\text{fm}
\end{equation}

\subsection{MOND Acceleration Scale ($a_0 = c H_\infty / (2\pi)$)}

The cosmological acceleration below which galactic rotation curves
flatten is derived from the asymptotic Hubble parameter:
\begin{equation}
    a_0 = \frac{c \cdot H_\infty}{2\pi}
    \approx 1.07 \times 10^{-10}\;\text{m/s}^2
\end{equation}
where $H_\infty = 28\pi m_e^3 c G / (\hbar^2 \alpha^2)$ (see
Eq.~\ref{eq:H_infinity}).  At this scale, the vacuum lattice's
mutual inductance begins to saturate, and the Axiom~4 saturation
term produces flat rotation curves identical to the Deep MOND
asymptote $g_{eff} \to \sqrt{g_N \cdot a_0}$ (see
\S\ref{sec:galactic_saturation}).

\subsection{Toroidal Halo Volume ($V_{halo} = 2$)}

The geometric volume of the 3D orthogonal tensor crossings of the
Borromean link is computed analytically from the signed intersection
integral of three great circles on $S^2$:
\begin{equation}
    V_{halo} = 2.0
\end{equation}
This is a purely topological constant; it sets the fraction of
the Skyrme soliton's energy that resides in the inter-braid halo
rather than the core.

\subsection{Protein Folding Barrier ($\beta = \ln(3) \times 3/7 = 0.471$)}

The folding timescale of a two-state protein is (Volume V, Chapter 5):
\begin{equation}
    \tau_\text{fold} = Q^2 \cdot N \cdot \tau_\text{water} \cdot \exp\!\bigl(\beta \cdot N \cdot \text{CO}\bigr)
\end{equation}
where $N$ is the chain length, CO is the relative contact order, and the barrier exponent combines three axiom-derived quantities:
\begin{align}
    \beta &= \ln(3) \times \frac{d}{n} = \ln(3) \times \frac{3}{7} = 0.471
\end{align}
\begin{itemize}
    \item $\ln(3)$: entropy cost of selecting 1 from 3 Ramachandran basins;
    \item $3/7 = d/n$: spatial compliance projection (same ratio as $\alpha_s = \alpha^{3/7}$);
    \item $Q^2 = 49$: double ring-down (propagation $\times$ evaluation);
    \item $\tau_\text{water} = 8.3$~ps: solvent rate-limiting.
\end{itemize}
Empirical barrier: $b_\text{emp} = 0.452$.  Derived: $\beta = 0.471$.  \textbf{Error: 4.1\%.}
Validated across 15 two-state folders: $R = 0.87$, 12/15 within $\pm$2 decades.

\subsection{FCC Packing Fraction ($\varphi = \pi\sqrt{2}/6$)}

The $K = 2G$ identity (Axiom~2) selects body-centred close-packing
of the lattice nodes.  The FCC packing fraction for equal spheres is:
\begin{equation}
    \varphi = \frac{\pi\sqrt{2}}{6} \approx 0.7405
\end{equation}
The complementary \textbf{void fraction} $1 - \varphi \approx 0.2595$ is the
interstitial space between packed nodes---geometrically inaccessible
to junction-based Op10 drain.

\paragraph{Dual role across regimes.}
At \textbf{nuclear scale} (Regime~I, $S \to 0$): the packing fraction
determines the saturated zone geometry that confines topological defects.
The crossing number $c$ partitions the Faddeev-Skyrme coupling via
$r_\text{opt} = \kappa/c$ --- the soliton cannot radiate past the packed
lattice boundary.

At \textbf{atomic scale} (Regime~II, $S \approx 1$): Op10 drain
$Y = c/(2\pi^2)$ acts on nodal phase space (fraction $\varphi$).
When $(1 - Y)^2 < 1 - \varphi$, the void fraction sets a floor:
\begin{equation}
    \text{IE} \geq E_\text{base} \times (1 - \varphi)
    \qquad \text{(void floor)}
\end{equation}
The threshold is $Y_c = 1 - \sqrt{1 - \varphi} \approx 0.49$.
This correction fixes F~($+0.2\%$) and Ne~($+0.9\%$) with zero
free parameters (see Volume IV, Chapter 16).

\subsection{Dielectric Yield Field ($E_{yield} = \sqrt{\alpha}\, E_S$)}

The Axiom~4 yield voltage $V_{yield} = \sqrt{\alpha}\, V_{snap} \approx 43.65$~kV applies across a single lattice node of spacing $\ell_{node} = \hbar/(m_e c) = 3.862 \times 10^{-13}$~m. The corresponding electric field threshold is:
\begin{equation}
    E_{yield} = \frac{V_{yield}}{\ell_{node}} = \frac{\sqrt{\alpha}\, m_e^2 c^3}{e\, \hbar} \approx 1.13 \times 10^{17}\;\text{V/m}
\end{equation}
This is exactly $\sqrt{\alpha}$ times the Schwinger critical field of QED ($E_S = m_e^2 c^3 / (e\hbar) \approx 1.32 \times 10^{18}$~V/m). The strongest sustained laboratory fields ($\sim 10^{10}$~V/m in ultrafast laser foci) are seven orders of magnitude below $E_{yield}$. Consequently, all macroscopic electromagnetic devices operate in \textbf{Regime~I} (linear vacuum, $\varepsilon_{eff} = \varepsilon_0$). Axiom~4 saturation is physically operative only at sub-femtometer separations---within particle cores, nuclear scattering events, and event horizons (see Volume IV, Chapter 2).

\subsection{Cooperative Lattice Amplification ($n_{coop} = 9$)}
\label{sec:n_coop_derivation}

The cooperative phase transition within hydrogen-bonded lattices is governed by a dimensionless amplification number $n_{coop}$ that converts single-bond energy into cooperative lattice yield strain. This constant admits two independent, convergent derivations.

\subsubsection*{Path~A: Tetrahedral Edge Count (Axiom~1)}

The hydrogen-bond network in water forms a tetrahedral coordination polyhedron. The cooperative unit is not the individual H-bond (vertex connection) but the \textbf{edge}---the pairwise coupling between two adjacent molecules. The tetrahedron has 4 vertices and \textbf{6 edges}.

Each edge is a one-dimensional coupling embedded in $D = 3$ spatial dimensions. The isotropic projection factor for a 2-body correlation in $D$ dimensions is $D/(D-1) = 3/2$. Therefore:
\begin{equation}
    n_{coop} = E_{\text{edges}} \times \frac{D}{D-1} = 6 \times \frac{3}{2} = 9
\end{equation}

\subsubsection*{Path~B: Compliance Mode Amplification (Axiom~2)}

The Poisson ratio $\nu_{vac} = 2/7$ implies \textbf{7 independent compliance modes} per lattice node ($\S$\ref{sec:g_star_derivation}). The bulk-shear coupling factor from the $K = 2G$ identity is $(1 + \nu_{vac}) = 9/7$. The cooperative amplification is then:
\begin{equation}
    n_{coop} = 7 \times (1 + \nu_{vac}) = 7 \times \frac{9}{7} = 9
\end{equation}

The exact algebraic identity $6 \times 3/2 = 7 \times 9/7 = 9$ confirms the structural consistency between the lattice geometry (Axiom~1) and the elastic compliance (Axiom~2).

\subsubsection*{Cross-Validation: Water +4$^\circ$C Density Anomaly}

The cooperative yield temperature $T_c$ is the point where the cooperative thermal strain amplitude $A = n_{coop} \cdot k_B T / E_{HB}$ reaches unity (i.e., the saturation operator yields):
\begin{equation}
    T_c = \frac{E_{HB}}{n_{coop} \cdot k_B}
        = \frac{0.2158\;\text{eV}}{9 \times 8.617 \times 10^{-5}\;\text{eV/K}}
        \approx 278.3\;\text{K}
        = +5.1^\circ\text{C}
\end{equation}

The measured anomalous density maximum of water occurs at $+4.0^\circ$C $= 277.15$~K.

\textbf{Error: 0.40\%.} Zero free parameters.

\subsubsection*{Biological Application: Cholesterol Phase Buffering}

In lipid bilayer membranes, the cooperative yield boundary at $T_c \approx 5^\circ$C governs the $V_I \to V_{II}$ gel-to-fluid transition. Cholesterol's rigid $sp^3$ 4-ring wedge raises the effective yield limit by the FCC packing fraction:
\begin{equation}
    A_{yield}^{\text{buffered}} = 1 + \varphi = 1 + \frac{\pi\sqrt{2}}{6} \approx 1.7405
\end{equation}
This pushes the catastrophic phase snap far outside the biological operating range ($-20^\circ$C to $+60^\circ$C), holding the membrane permanently at the $K = 2G$ structural yield threshold---pliant enough for wave propagation, rigid enough for allosteric state-switching in embedded protein circuits (Volume V, Chapter 2, $\S$\ref{sec:membrane_phase_buffering}).

\section{$g_* = 7^3/4 = 85.75$: Lattice Mode Count}
\label{sec:g_star_derivation}

The effective number of relativistic degrees of freedom $g_*$
appears in the baryon asymmetry formula.  The Standard Model
counts $g_{*,SM} = 106.75$ by enumerating all known particles.
AVE derives this constant from lattice geometry.

\subsection{Derivation}

The Poisson ratio $\nu_{vac} = 2/7$ (derived in Book~1) reveals
that each lattice node has \textbf{7 independent compliance modes}:
the denominator of $\nu$.  Physically, these are the 7 independent
small-deformation degrees of freedom of the LC network at each
node: 3 translational, 3 rotational, and 1 volumetric.

In three dimensions, the total mode count is $7^3 = 343$ (7 per
spatial direction, cubed for the full 3D Brillouin zone).  The
K4 unit cell contains 4 nodes, so the effective DOF per cell is:
\begin{equation}
    \boxed{g_* = \frac{7^3}{4} = \frac{343}{4} = 85.75}
\end{equation}

\subsection{Verification}

Substituting into the baryon asymmetry formula with all other
factors derived from AVE constants:
\begin{equation}
    \eta = \frac{(\pi/\kappa_{FS}) \cdot \alpha_W^4 \cdot (28/79)}{7^3/4}
    = 6.08 \times 10^{-10}
\end{equation}

Observed: $\eta_{obs} = 6.1 \times 10^{-10}$.
\textbf{Error: 0.38\%.}  Using the SM value $g_* = 106.75$
gives 20\% error.

\section{$\alpha_s = \alpha^{3/7} \approx 0.121$: Strong Coupling}
\label{sec:alpha_s_derivation}

The strong coupling constant $\alpha_s$ at the $Z$-pole ($M_Z$)
is measured as $0.1179 \pm 0.0010$ (PDG~\cite{pdg2022}).  AVE derives it from
the lattice compliance structure.

\subsection{Derivation}

The EM coupling $\alpha$ operates on the full 7-mode compliance
manifold of each lattice node (the same 7 that gives $\nu_{vac} = 2/7$
and $g_* = 7^3/4$).  The strong coupling is the \textbf{spatial
projection}: $\alpha$ restricted to the 3 translational modes
out of 7 total.

The projected coupling scales as a fractional power:
\begin{equation}
    \boxed{\alpha_s = \alpha^{d/n} = \alpha^{3/7} \approx 0.1214}
\end{equation}
where $d = 3$ (spatial dimensions) and $n = 7$ (compliance modes
per node, from $\nu_{vac} = 2/7$).

\textbf{Error: 2.97\%.}  Zero free parameters.

\subsection{Physical Interpretation}

$\alpha$ is the coupling strength for ALL 7 lattice modes (3
translational + 3 rotational + 1 volumetric).  When the
interaction is confined to spatial translations only (as in the
strong force between quarks), the effective coupling is the
$3/7$-th power of $\alpha$.  This is a dimensional projection,
not a fit.


\section{$\lambda_H = 1/8$: Higgs Quartic Coupling}
\label{sec:lambda_higgs_derivation}

The Standard Model Higgs quartic self-coupling $\lambda$ is
measured as $\approx 0.129$.  AVE derives it from the K4 unit
cell topology.

\subsection{Derivation}

The Higgs boson is the radial \textbf{breathing mode} of the
K4 unit cell: a symmetric compression/expansion of all 4
nodes simultaneously.  The quartic restoring force of this
mode is shared equally across $N_{K4} = 4$ nodes:
\begin{equation}
    \boxed{\lambda_H = \frac{1}{2\,N_{K4}} = \frac{1}{8} = 0.125}
\end{equation}

The Higgs mass follows from the standard relation $m_H = \sqrt{2\lambda}\,v$:
\begin{equation}
    m_H = v\,\sqrt{\frac{2}{2\,N_{K4}}} = \frac{v}{\sqrt{N_{K4}}}
    = \frac{v}{2} \approx 124{,}400\;\text{MeV}
\end{equation}

Observed: $m_H = 125{,}100$~MeV.  \textbf{Error: 0.55\%.}

\subsection{Why $N_{K4} = 4$?}

The K4 complete graph is the unique 4-vertex graph where every
node connects to every other node.  It is the minimal unit cell
of the SRS lattice (the chiral rod packing that defines the vacuum).
The number 4 is not arbitrary: it is the minimum number of vertices
required for a 3D chiral simplex.


\section{Conclusion: The Death of Numerology}

The Golden Ratio is the icosahedron. The Fibonacci sequence is an icosahedral approximation algorithm. $\pi$ is Gauss's law. The magic numbers are geometric shell closures. $g_* = 7^3/4$ is the lattice mode count. $\alpha_s = \alpha^{3/7}$ is the spatial compliance projection. $\lambda_H = 1/8$ is the K4 breathing mode. $\beta_\text{fold} = \ln(3) \times 3/7 = 0.471$ is the protein folding barrier exponent, using the \emph{same} $d/n$ ratio as $\alpha_s$. $n_{coop} = 9 = 6 \times 3/2 = 7 \times 9/7$ is the cooperative lattice amplification that derives water's $+4^\circ$C density anomaly to $0.40\%$. $T_m = \hbar\sqrt{k_{hb}/m_H}/k_B = 279.5$~K is the proton transfer eigenmode of the O--H$\cdots$O bridge, deriving the melting point of water to $+2.3\%$ with zero spectroscopic inputs.

Every instance of a ``mystical'' mathematical constant appearing in nature reduces, upon derivation, to a packing theorem. The packing theorems themselves are consequences of minimizing mutual reactive impedance between discrete topological defects in a bounded 3D medium. No design principle, cosmic harmony, or transcendent mathematics is required---only geometry and the requirement that coupled LC oscillators minimize total strain.

Applied Vacuum Engineering does not discover new mathematical constants. It derives the ones that numerology could only observe.
